\documentclass[10pt,a4paper]{amsart}
\usepackage{amsmath,amssymb,amsthm,graphicx}
% Retained minimal preamble for the Stacks Project chapter more-morphisms.tex (More on Morphisms).
% Verbatim excerpt of the chapter preamble: Xy-pic setup (native preamble lines 14-18) and the native
% theorem declarations with the mathematical macros (native preamble lines 154-217). Copied mechanically;
% nothing added, removed or redefined. The wrapper class and the amsmath/amssymb/amsthm/graphicx base
% packages are supplied by the assembler, so the class/externaldocument part of the native preamble is omitted.
% For commutative diagrams we use Xy-pic
\usepackage[all]{xy}

% We use 2cell for 2-commutative diagrams.
\xyoption{2cell}
% Theorem environments.
%
\theoremstyle{plain}
\newtheorem{theorem}[subsection]{Theorem}
\newtheorem{proposition}[subsection]{Proposition}
\newtheorem{lemma}[subsection]{Lemma}

\theoremstyle{definition}
\newtheorem{definition}[subsection]{Definition}
\newtheorem{example}[subsection]{Example}
\newtheorem{exercise}[subsection]{Exercise}
\newtheorem{situation}[subsection]{Situation}

\theoremstyle{remark}
\newtheorem{remark}[subsection]{Remark}
\newtheorem{remarks}[subsection]{Remarks}

\numberwithin{equation}{subsection}

% Macros
%
\def\lim{\mathop{\mathrm{lim}}\nolimits}
\def\colim{\mathop{\mathrm{colim}}\nolimits}
\def\Spec{\mathop{\mathrm{Spec}}}
\def\Hom{\mathop{\mathrm{Hom}}\nolimits}
\def\Ext{\mathop{\mathrm{Ext}}\nolimits}
\def\SheafHom{\mathop{\mathcal{H}\!\mathit{om}}\nolimits}
\def\SheafExt{\mathop{\mathcal{E}\!\mathit{xt}}\nolimits}
\def\Sch{\mathit{Sch}}
\def\Mor{\mathop{\mathrm{Mor}}\nolimits}
\def\Ob{\mathop{\mathrm{Ob}}\nolimits}
\def\Sh{\mathop{\mathit{Sh}}\nolimits}
\def\NL{\mathop{N\!L}\nolimits}
\def\CH{\mathop{\mathrm{CH}}\nolimits}
\def\proetale{{pro\text{-}\acute{e}tale}}
\def\etale{{\acute{e}tale}}
\def\QCoh{\mathit{QCoh}}
\def\Ker{\mathop{\mathrm{Ker}}}
\def\Im{\mathop{\mathrm{Im}}}
\def\Coker{\mathop{\mathrm{Coker}}}
\def\Coim{\mathop{\mathrm{Coim}}}

% Boxtimes
%
\DeclareMathSymbol{\boxtimes}{\mathbin}{AMSa}{"02}

%
% Macros for moduli stacks/spaces
%
\def\QCohstack{\mathcal{QC}\!\mathit{oh}}
\def\Cohstack{\mathcal{C}\!\mathit{oh}}
\def\Spacesstack{\mathcal{S}\!\mathit{paces}}
\def\Quotfunctor{\mathrm{Quot}}
\def\Hilbfunctor{\mathrm{Hilb}}
\def\Curvesstack{\mathcal{C}\!\mathit{urves}}
\def\Polarizedstack{\mathcal{P}\!\mathit{olarized}}
\def\Complexesstack{\mathcal{C}\!\mathit{omplexes}}
% \Pic is the operator that assigns to X its picard group, usage \Pic(X)
% \Picardstack_{X/B} denotes the Picard stack of X over B
% \Picardfunctor_{X/B} denotes the Picard functor of X over B
\def\Pic{\mathop{\mathrm{Pic}}\nolimits}
\def\Picardstack{\mathcal{P}\!\mathit{ic}}
\def\Picardfunctor{\mathrm{Pic}}
\def\Deformationcategory{\mathcal{D}\!\mathit{ef}}
\usepackage{cleveref}
\crefname{originallemma}{original lemma}{original lemmas}
\Crefname{originallemma}{Original lemma}{Original lemmas}
\makeatletter
\newlabel{algebra-lemma-NAK}{{00DV}{source}{}{sourcepointer-1}{}}
\newlabel{algebra-lemma-NAK@cref}{{[originallemma][0][]00DV}{[1][0][]source}{}{}{}}
\makeatother
\crefname{originallemma}{original lemma}{original lemmas}
\Crefname{originallemma}{Original lemma}{Original lemmas}
\makeatletter
\newlabel{definition-thickening}{{definition-thickening}{source}{}{sourcepointer-2}{}}
\newlabel{definition-thickening@cref}{{[originallemma][0][]definition-thickening}{[1][0][]source}{}{}{}}
\makeatother
\crefname{originallemma}{original lemma}{original lemmas}
\Crefname{originallemma}{Original lemma}{Original lemmas}
\makeatletter
\newlabel{lemma-deform}{{lemma-deform}{source}{}{sourcepointer-3}{}}
\newlabel{lemma-deform@cref}{{[originallemma][0][]lemma-deform}{[1][0][]source}{}{}{}}
\makeatother
\crefname{originallemma}{original lemma}{original lemmas}
\Crefname{originallemma}{Original lemma}{Original lemmas}
\makeatletter
\newlabel{lemma-flatness-morphism-thickenings}{{lemma-flatness-morphism-thickenings}{source}{}{sourcepointer-4}{}}
\newlabel{lemma-flatness-morphism-thickenings@cref}{{[originallemma][0][]lemma-flatness-morphism-thickenings}{[1][0][]source}{}{}{}}
\makeatother
\crefname{originallemma}{original lemma}{original lemmas}
\Crefname{originallemma}{Original lemma}{Original lemmas}
\makeatletter
\newlabel{lemma-thicken-property-morphisms}{{lemma-thicken-property-morphisms}{source}{}{sourcepointer-5}{}}
\newlabel{lemma-thicken-property-morphisms@cref}{{[originallemma][0][]lemma-thicken-property-morphisms}{[1][0][]source}{}{}{}}
\makeatother
\crefname{originallemma}{original lemma}{original lemmas}
\Crefname{originallemma}{Original lemma}{Original lemmas}
\makeatletter
\newlabel{lemma-thicken-property-morphisms-cartesian}{{lemma-thicken-property-morphisms-cartesian}{source}{}{sourcepointer-6}{}}
\newlabel{lemma-thicken-property-morphisms-cartesian@cref}{{[originallemma][0][]lemma-thicken-property-morphisms-cartesian}{[1][0][]source}{}{}{}}
\makeatother
\crefname{originallemma}{original lemma}{original lemmas}
\Crefname{originallemma}{Original lemma}{Original lemmas}
\makeatletter
\newlabel{lemma-thickening-affine-scheme}{{lemma-thickening-affine-scheme}{source}{}{sourcepointer-7}{}}
\newlabel{lemma-thickening-affine-scheme@cref}{{[originallemma][0][]lemma-thickening-affine-scheme}{[1][0][]source}{}{}{}}
\makeatother
\crefname{originallemma}{original lemma}{original lemmas}
\Crefname{originallemma}{Original lemma}{Original lemmas}
\makeatletter
\newlabel{morphisms-lemma-base-change-affine}{{01SD}{source}{}{sourcepointer-8}{}}
\newlabel{morphisms-lemma-base-change-affine@cref}{{[originallemma][0][]01SD}{[1][0][]source}{}{}{}}
\makeatother
\crefname{originallemma}{original lemma}{original lemmas}
\Crefname{originallemma}{Original lemma}{Original lemmas}
\makeatletter
\newlabel{morphisms-lemma-base-change-etale}{{02GO}{source}{}{sourcepointer-9}{}}
\newlabel{morphisms-lemma-base-change-etale@cref}{{[originallemma][0][]02GO}{[1][0][]source}{}{}{}}
\makeatother
\crefname{originallemma}{original lemma}{original lemmas}
\Crefname{originallemma}{Original lemma}{Original lemmas}
\makeatletter
\newlabel{morphisms-lemma-base-change-finite}{{01WL}{source}{}{sourcepointer-10}{}}
\newlabel{morphisms-lemma-base-change-finite@cref}{{[originallemma][0][]01WL}{[1][0][]source}{}{}{}}
\makeatother
\crefname{originallemma}{original lemma}{original lemmas}
\Crefname{originallemma}{Original lemma}{Original lemmas}
\makeatletter
\newlabel{morphisms-lemma-base-change-finite-locally-free}{{02KD}{source}{}{sourcepointer-11}{}}
\newlabel{morphisms-lemma-base-change-finite-locally-free@cref}{{[originallemma][0][]02KD}{[1][0][]source}{}{}{}}
\makeatother
\crefname{originallemma}{original lemma}{original lemmas}
\Crefname{originallemma}{Original lemma}{Original lemmas}
\makeatletter
\newlabel{morphisms-lemma-base-change-finite-presentation}{{01TS}{source}{}{sourcepointer-12}{}}
\newlabel{morphisms-lemma-base-change-finite-presentation@cref}{{[originallemma][0][]01TS}{[1][0][]source}{}{}{}}
\makeatother
\crefname{originallemma}{original lemma}{original lemmas}
\Crefname{originallemma}{Original lemma}{Original lemmas}
\makeatletter
\newlabel{morphisms-lemma-base-change-finite-type}{{01T4}{source}{}{sourcepointer-13}{}}
\newlabel{morphisms-lemma-base-change-finite-type@cref}{{[originallemma][0][]01T4}{[1][0][]source}{}{}{}}
\makeatother
\crefname{originallemma}{original lemma}{original lemmas}
\Crefname{originallemma}{Original lemma}{Original lemmas}
\makeatletter
\newlabel{morphisms-lemma-base-change-proper}{{01W4}{source}{}{sourcepointer-14}{}}
\newlabel{morphisms-lemma-base-change-proper@cref}{{[originallemma][0][]01W4}{[1][0][]source}{}{}{}}
\makeatother
\crefname{originallemma}{original lemma}{original lemmas}
\Crefname{originallemma}{Original lemma}{Original lemmas}
\makeatletter
\newlabel{morphisms-lemma-base-change-quasi-finite}{{01TM}{source}{}{sourcepointer-15}{}}
\newlabel{morphisms-lemma-base-change-quasi-finite@cref}{{[originallemma][0][]01TM}{[1][0][]source}{}{}{}}
\makeatother
\crefname{originallemma}{original lemma}{original lemmas}
\Crefname{originallemma}{Original lemma}{Original lemmas}
\makeatletter
\newlabel{morphisms-lemma-base-change-relative-dimension-d}{{02NK}{source}{}{sourcepointer-16}{}}
\newlabel{morphisms-lemma-base-change-relative-dimension-d@cref}{{[originallemma][0][]02NK}{[1][0][]source}{}{}{}}
\makeatother
\crefname{originallemma}{original lemma}{original lemmas}
\Crefname{originallemma}{Original lemma}{Original lemmas}
\makeatletter
\newlabel{morphisms-lemma-base-change-smooth}{{01VB}{source}{}{sourcepointer-17}{}}
\newlabel{morphisms-lemma-base-change-smooth@cref}{{[originallemma][0][]01VB}{[1][0][]source}{}{}{}}
\makeatother
\crefname{originallemma}{original lemma}{original lemmas}
\Crefname{originallemma}{Original lemma}{Original lemmas}
\makeatletter
\newlabel{morphisms-lemma-base-change-surjective}{{01S1}{source}{}{sourcepointer-18}{}}
\newlabel{morphisms-lemma-base-change-surjective@cref}{{[originallemma][0][]01S1}{[1][0][]source}{}{}{}}
\makeatother
\crefname{originallemma}{original lemma}{original lemmas}
\Crefname{originallemma}{Original lemma}{Original lemmas}
\makeatletter
\newlabel{morphisms-lemma-base-change-syntomic}{{01UI}{source}{}{sourcepointer-19}{}}
\newlabel{morphisms-lemma-base-change-syntomic@cref}{{[originallemma][0][]01UI}{[1][0][]source}{}{}{}}
\makeatother
\crefname{originallemma}{original lemma}{original lemmas}
\Crefname{originallemma}{Original lemma}{Original lemmas}
\makeatletter
\newlabel{morphisms-lemma-base-change-universally-injective}{{0472}{source}{}{sourcepointer-20}{}}
\newlabel{morphisms-lemma-base-change-universally-injective@cref}{{[originallemma][0][]0472}{[1][0][]source}{}{}{}}
\makeatother
\crefname{originallemma}{original lemma}{original lemmas}
\Crefname{originallemma}{Original lemma}{Original lemmas}
\makeatletter
\newlabel{morphisms-lemma-base-change-unramified}{{02GA}{source}{}{sourcepointer-21}{}}
\newlabel{morphisms-lemma-base-change-unramified@cref}{{[originallemma][0][]02GA}{[1][0][]source}{}{}{}}
\makeatother
\crefname{originallemma}{original lemma}{original lemmas}
\Crefname{originallemma}{Original lemma}{Original lemmas}
\makeatletter
\newlabel{morphisms-lemma-etale-flat-etale-fibres}{{02GM}{source}{}{sourcepointer-22}{}}
\newlabel{morphisms-lemma-etale-flat-etale-fibres@cref}{{[originallemma][0][]02GM}{[1][0][]source}{}{}{}}
\makeatother
\crefname{originallemma}{original lemma}{original lemmas}
\Crefname{originallemma}{Original lemma}{Original lemmas}
\makeatletter
\newlabel{morphisms-lemma-finite-flat}{{02KB}{source}{}{sourcepointer-23}{}}
\newlabel{morphisms-lemma-finite-flat@cref}{{[originallemma][0][]02KB}{[1][0][]source}{}{}{}}
\makeatother
\crefname{originallemma}{original lemma}{original lemmas}
\Crefname{originallemma}{Original lemma}{Original lemmas}
\makeatletter
\newlabel{morphisms-lemma-finite-integral}{{01WJ}{source}{}{sourcepointer-24}{}}
\newlabel{morphisms-lemma-finite-integral@cref}{{[originallemma][0][]01WJ}{[1][0][]source}{}{}{}}
\makeatother
\crefname{originallemma}{original lemma}{original lemmas}
\Crefname{originallemma}{Original lemma}{Original lemmas}
\makeatletter
\newlabel{morphisms-lemma-integral-universally-closed}{{01WM}{source}{}{sourcepointer-25}{}}
\newlabel{morphisms-lemma-integral-universally-closed@cref}{{[originallemma][0][]01WM}{[1][0][]source}{}{}{}}
\makeatother
\crefname{originallemma}{original lemma}{original lemmas}
\Crefname{originallemma}{Original lemma}{Original lemmas}
\makeatletter
\newlabel{morphisms-lemma-quasi-finite-at-point-characterize}{{01TH}{source}{}{sourcepointer-26}{}}
\newlabel{morphisms-lemma-quasi-finite-at-point-characterize@cref}{{[originallemma][0][]01TH}{[1][0][]source}{}{}{}}
\makeatother
\crefname{originallemma}{original lemma}{original lemmas}
\Crefname{originallemma}{Original lemma}{Original lemmas}
\makeatletter
\newlabel{morphisms-lemma-smooth-flat-smooth-fibres}{{01V8}{source}{}{sourcepointer-27}{}}
\newlabel{morphisms-lemma-smooth-flat-smooth-fibres@cref}{{[originallemma][0][]01V8}{[1][0][]source}{}{}{}}
\makeatother
\crefname{originallemma}{original lemma}{original lemmas}
\Crefname{originallemma}{Original lemma}{Original lemmas}
\makeatletter
\newlabel{morphisms-lemma-syntomic-flat-fibres}{{01UF}{source}{}{sourcepointer-28}{}}
\newlabel{morphisms-lemma-syntomic-flat-fibres@cref}{{[originallemma][0][]01UF}{[1][0][]source}{}{}{}}
\makeatother
\crefname{originallemma}{original lemma}{original lemmas}
\Crefname{originallemma}{Original lemma}{Original lemmas}
\makeatletter
\newlabel{morphisms-lemma-unramified-etale-fibres}{{02G8}{source}{}{sourcepointer-29}{}}
\newlabel{morphisms-lemma-unramified-etale-fibres@cref}{{[originallemma][0][]02G8}{[1][0][]source}{}{}{}}
\makeatother
\crefname{originallemma}{original lemma}{original lemmas}
\Crefname{originallemma}{Original lemma}{Original lemmas}
\makeatletter
\newlabel{schemes-lemma-base-change-immersion}{{01JY}{source}{}{sourcepointer-30}{}}
\newlabel{schemes-lemma-base-change-immersion@cref}{{[originallemma][0][]01JY}{[1][0][]source}{}{}{}}
\makeatother
\crefname{originallemma}{original lemma}{original lemmas}
\Crefname{originallemma}{Original lemma}{Original lemmas}
\makeatletter
\newlabel{schemes-lemma-base-change-monomorphism}{{02YC}{source}{}{sourcepointer-31}{}}
\newlabel{schemes-lemma-base-change-monomorphism@cref}{{[originallemma][0][]02YC}{[1][0][]source}{}{}{}}
\makeatother
\crefname{originallemma}{original lemma}{original lemmas}
\Crefname{originallemma}{Original lemma}{Original lemmas}
\makeatletter
\newlabel{schemes-lemma-quasi-compact-preserved-base-change}{{01K5}{source}{}{sourcepointer-32}{}}
\newlabel{schemes-lemma-quasi-compact-preserved-base-change@cref}{{[originallemma][0][]01K5}{[1][0][]source}{}{}{}}
\makeatother
\crefname{originallemma}{original lemma}{original lemmas}
\Crefname{originallemma}{Original lemma}{Original lemmas}
\makeatletter
\newlabel{schemes-lemma-separated-permanence}{{01KU}{source}{}{sourcepointer-33}{}}
\newlabel{schemes-lemma-separated-permanence@cref}{{[originallemma][0][]01KU}{[1][0][]source}{}{}{}}
\makeatother
\title{Twenty-three local properties of a morphism are insensitive to extending a finite order thickening of the target}
\author{The Stacks Project Authors}
\date{stacks-project:more-morphisms.tex:local-snapshot-20261011; GNU FDL 1.2 or later}
\begin{document}
\maketitle
\noindent\textit{Editorial note.} Statement span is lines 2266-2334 including the initial xymatrix and the twenty-three enumerated items; proof span is lines 2336-2560, the complete author proof of 225 lines. This is the largest proof collected in this lane so far.\par\medskip
\noindent\textit{Editorial note.} The statement's own twenty-third item reads 'add more here', the author's placeholder; it is retained verbatim because it belongs to the collected statement environment.\par\medskip
\noindent\textit{Editorial note.} The proof's forward direction is the author's own citation list of base-change stability results; the converse direction and the finite locally free endgame are fully carried out in the retained text, which is what the completeness rule requires.\par\medskip
\noindent\textit{Editorial note.} This card carries the largest navigation-stub set in this lane (33 cited-but-not-included labels); all are prerequisites and are marked non-load-bearing.\par\medskip
\noindent\textit{Editorial note.} extra\_packages is empty on purpose: the retained preamble already loads xy with options and a second declaration clashes.\par\medskip
\noindent\textit{Allowed standard prerequisites (editorial).}\begin{itemize}
\item Explicit prerequisite boundary. Stability of the listed properties under base change, finite order thickenings, and the characterisations of properness, finiteness and finite local freeness. The converse direction of the equivalence and the endgame are retained; the forward direction is the author's own citation list.
\end{itemize}\medskip
\ifdefined\hypertarget\hypertarget{sourcepointer-1}{}\fi
\noindent\textit{Source reference.} Original lemma 00DV: sources/stacks-more-morphisms/pinned-original-source.tex, source lines 2560--2560. Referenced by the retained statement or proof but not part of the collected spans; supplied only as a navigation pointer. The argument does not depend on its content.\par\medskip
\ifdefined\hypertarget\hypertarget{sourcepointer-2}{}\fi
\noindent\textit{Source reference.} Original lemma definition-thickening: sources/stacks-more-morphisms/pinned-original-source.tex, source lines 36--63. Referenced by the retained statement or proof but not part of the collected spans; supplied only as a navigation pointer. The argument does not depend on its content.\par\medskip
\ifdefined\hypertarget\hypertarget{sourcepointer-3}{}\fi
\noindent\textit{Source reference.} Original lemma lemma-deform: sources/stacks-more-morphisms/pinned-original-source.tex, source lines 2182--2190. Referenced by the retained statement or proof but not part of the collected spans; supplied only as a navigation pointer. The argument does not depend on its content.\par\medskip
\ifdefined\hypertarget\hypertarget{sourcepointer-4}{}\fi
\noindent\textit{Source reference.} Original lemma lemma-flatness-morphism-thickenings: sources/stacks-more-morphisms/pinned-original-source.tex, source lines 2238--2252. Referenced by the retained statement or proof but not part of the collected spans; supplied only as a navigation pointer. The argument does not depend on its content.\par\medskip
\ifdefined\hypertarget\hypertarget{sourcepointer-5}{}\fi
\noindent\textit{Source reference.} Original lemma lemma-thicken-property-morphisms: sources/stacks-more-morphisms/pinned-original-source.tex, source lines 270--284. Referenced by the retained statement or proof but not part of the collected spans; supplied only as a navigation pointer. The argument does not depend on its content.\par\medskip
\ifdefined\hypertarget\hypertarget{sourcepointer-6}{}\fi
\noindent\textit{Source reference.} Original lemma lemma-thicken-property-morphisms-cartesian: sources/stacks-more-morphisms/pinned-original-source.tex, source lines 335--354. Referenced by the retained statement or proof but not part of the collected spans; supplied only as a navigation pointer. The argument does not depend on its content.\par\medskip
\ifdefined\hypertarget\hypertarget{sourcepointer-7}{}\fi
\noindent\textit{Source reference.} Original lemma lemma-thickening-affine-scheme: sources/stacks-more-morphisms/pinned-original-source.tex, source lines 157--160. Referenced by the retained statement or proof but not part of the collected spans; supplied only as a navigation pointer. The argument does not depend on its content.\par\medskip
\ifdefined\hypertarget\hypertarget{sourcepointer-8}{}\fi
\noindent\textit{Source reference.} Original lemma 01SD: sources/stacks-more-morphisms/pinned-original-source.tex, source lines 2560--2560. Referenced by the retained statement or proof but not part of the collected spans; supplied only as a navigation pointer. The argument does not depend on its content.\par\medskip
\ifdefined\hypertarget\hypertarget{sourcepointer-9}{}\fi
\noindent\textit{Source reference.} Original lemma 02GO: sources/stacks-more-morphisms/pinned-original-source.tex, source lines 2560--2560. Referenced by the retained statement or proof but not part of the collected spans; supplied only as a navigation pointer. The argument does not depend on its content.\par\medskip
\ifdefined\hypertarget\hypertarget{sourcepointer-10}{}\fi
\noindent\textit{Source reference.} Original lemma 01WL: sources/stacks-more-morphisms/pinned-original-source.tex, source lines 2560--2560. Referenced by the retained statement or proof but not part of the collected spans; supplied only as a navigation pointer. The argument does not depend on its content.\par\medskip
\ifdefined\hypertarget\hypertarget{sourcepointer-11}{}\fi
\noindent\textit{Source reference.} Original lemma 02KD: sources/stacks-more-morphisms/pinned-original-source.tex, source lines 2560--2560. Referenced by the retained statement or proof but not part of the collected spans; supplied only as a navigation pointer. The argument does not depend on its content.\par\medskip
\ifdefined\hypertarget\hypertarget{sourcepointer-12}{}\fi
\noindent\textit{Source reference.} Original lemma 01TS: sources/stacks-more-morphisms/pinned-original-source.tex, source lines 2560--2560. Referenced by the retained statement or proof but not part of the collected spans; supplied only as a navigation pointer. The argument does not depend on its content.\par\medskip
\ifdefined\hypertarget\hypertarget{sourcepointer-13}{}\fi
\noindent\textit{Source reference.} Original lemma 01T4: sources/stacks-more-morphisms/pinned-original-source.tex, source lines 2560--2560. Referenced by the retained statement or proof but not part of the collected spans; supplied only as a navigation pointer. The argument does not depend on its content.\par\medskip
\ifdefined\hypertarget\hypertarget{sourcepointer-14}{}\fi
\noindent\textit{Source reference.} Original lemma 01W4: sources/stacks-more-morphisms/pinned-original-source.tex, source lines 2560--2560. Referenced by the retained statement or proof but not part of the collected spans; supplied only as a navigation pointer. The argument does not depend on its content.\par\medskip
\ifdefined\hypertarget\hypertarget{sourcepointer-15}{}\fi
\noindent\textit{Source reference.} Original lemma 01TM: sources/stacks-more-morphisms/pinned-original-source.tex, source lines 2560--2560. Referenced by the retained statement or proof but not part of the collected spans; supplied only as a navigation pointer. The argument does not depend on its content.\par\medskip
\ifdefined\hypertarget\hypertarget{sourcepointer-16}{}\fi
\noindent\textit{Source reference.} Original lemma 02NK: sources/stacks-more-morphisms/pinned-original-source.tex, source lines 2560--2560. Referenced by the retained statement or proof but not part of the collected spans; supplied only as a navigation pointer. The argument does not depend on its content.\par\medskip
\ifdefined\hypertarget\hypertarget{sourcepointer-17}{}\fi
\noindent\textit{Source reference.} Original lemma 01VB: sources/stacks-more-morphisms/pinned-original-source.tex, source lines 2560--2560. Referenced by the retained statement or proof but not part of the collected spans; supplied only as a navigation pointer. The argument does not depend on its content.\par\medskip
\ifdefined\hypertarget\hypertarget{sourcepointer-18}{}\fi
\noindent\textit{Source reference.} Original lemma 01S1: sources/stacks-more-morphisms/pinned-original-source.tex, source lines 2560--2560. Referenced by the retained statement or proof but not part of the collected spans; supplied only as a navigation pointer. The argument does not depend on its content.\par\medskip
\ifdefined\hypertarget\hypertarget{sourcepointer-19}{}\fi
\noindent\textit{Source reference.} Original lemma 01UI: sources/stacks-more-morphisms/pinned-original-source.tex, source lines 2560--2560. Referenced by the retained statement or proof but not part of the collected spans; supplied only as a navigation pointer. The argument does not depend on its content.\par\medskip
\ifdefined\hypertarget\hypertarget{sourcepointer-20}{}\fi
\noindent\textit{Source reference.} Original lemma 0472: sources/stacks-more-morphisms/pinned-original-source.tex, source lines 2560--2560. Referenced by the retained statement or proof but not part of the collected spans; supplied only as a navigation pointer. The argument does not depend on its content.\par\medskip
\ifdefined\hypertarget\hypertarget{sourcepointer-21}{}\fi
\noindent\textit{Source reference.} Original lemma 02GA: sources/stacks-more-morphisms/pinned-original-source.tex, source lines 2560--2560. Referenced by the retained statement or proof but not part of the collected spans; supplied only as a navigation pointer. The argument does not depend on its content.\par\medskip
\ifdefined\hypertarget\hypertarget{sourcepointer-22}{}\fi
\noindent\textit{Source reference.} Original lemma 02GM: sources/stacks-more-morphisms/pinned-original-source.tex, source lines 2560--2560. Referenced by the retained statement or proof but not part of the collected spans; supplied only as a navigation pointer. The argument does not depend on its content.\par\medskip
\ifdefined\hypertarget\hypertarget{sourcepointer-23}{}\fi
\noindent\textit{Source reference.} Original lemma 02KB: sources/stacks-more-morphisms/pinned-original-source.tex, source lines 2560--2560. Referenced by the retained statement or proof but not part of the collected spans; supplied only as a navigation pointer. The argument does not depend on its content.\par\medskip
\ifdefined\hypertarget\hypertarget{sourcepointer-24}{}\fi
\noindent\textit{Source reference.} Original lemma 01WJ: sources/stacks-more-morphisms/pinned-original-source.tex, source lines 2560--2560. Referenced by the retained statement or proof but not part of the collected spans; supplied only as a navigation pointer. The argument does not depend on its content.\par\medskip
\ifdefined\hypertarget\hypertarget{sourcepointer-25}{}\fi
\noindent\textit{Source reference.} Original lemma 01WM: sources/stacks-more-morphisms/pinned-original-source.tex, source lines 2560--2560. Referenced by the retained statement or proof but not part of the collected spans; supplied only as a navigation pointer. The argument does not depend on its content.\par\medskip
\ifdefined\hypertarget\hypertarget{sourcepointer-26}{}\fi
\noindent\textit{Source reference.} Original lemma 01TH: sources/stacks-more-morphisms/pinned-original-source.tex, source lines 2560--2560. Referenced by the retained statement or proof but not part of the collected spans; supplied only as a navigation pointer. The argument does not depend on its content.\par\medskip
\ifdefined\hypertarget\hypertarget{sourcepointer-27}{}\fi
\noindent\textit{Source reference.} Original lemma 01V8: sources/stacks-more-morphisms/pinned-original-source.tex, source lines 2560--2560. Referenced by the retained statement or proof but not part of the collected spans; supplied only as a navigation pointer. The argument does not depend on its content.\par\medskip
\ifdefined\hypertarget\hypertarget{sourcepointer-28}{}\fi
\noindent\textit{Source reference.} Original lemma 01UF: sources/stacks-more-morphisms/pinned-original-source.tex, source lines 2560--2560. Referenced by the retained statement or proof but not part of the collected spans; supplied only as a navigation pointer. The argument does not depend on its content.\par\medskip
\ifdefined\hypertarget\hypertarget{sourcepointer-29}{}\fi
\noindent\textit{Source reference.} Original lemma 02G8: sources/stacks-more-morphisms/pinned-original-source.tex, source lines 2560--2560. Referenced by the retained statement or proof but not part of the collected spans; supplied only as a navigation pointer. The argument does not depend on its content.\par\medskip
\ifdefined\hypertarget\hypertarget{sourcepointer-30}{}\fi
\noindent\textit{Source reference.} Original lemma 01JY: sources/stacks-more-morphisms/pinned-original-source.tex, source lines 2560--2560. Referenced by the retained statement or proof but not part of the collected spans; supplied only as a navigation pointer. The argument does not depend on its content.\par\medskip
\ifdefined\hypertarget\hypertarget{sourcepointer-31}{}\fi
\noindent\textit{Source reference.} Original lemma 02YC: sources/stacks-more-morphisms/pinned-original-source.tex, source lines 2560--2560. Referenced by the retained statement or proof but not part of the collected spans; supplied only as a navigation pointer. The argument does not depend on its content.\par\medskip
\ifdefined\hypertarget\hypertarget{sourcepointer-32}{}\fi
\noindent\textit{Source reference.} Original lemma 01K5: sources/stacks-more-morphisms/pinned-original-source.tex, source lines 2560--2560. Referenced by the retained statement or proof but not part of the collected spans; supplied only as a navigation pointer. The argument does not depend on its content.\par\medskip
\ifdefined\hypertarget\hypertarget{sourcepointer-33}{}\fi
\noindent\textit{Source reference.} Original lemma 01KU: sources/stacks-more-morphisms/pinned-original-source.tex, source lines 2560--2560. Referenced by the retained statement or proof but not part of the collected spans; supplied only as a navigation pointer. The argument does not depend on its content.\par\medskip
\begin{lemma}
\label{lemma-deform-property}
Consider a commutative diagram
$$
\xymatrix{
(X \subset X') \ar[rr]_{(f, f')} \ar[rd] & & (Y \subset Y') \ar[ld] \\
& (S \subset S')
}
$$
of thickenings. Assume $S \subset S'$ is a finite order thickening,
$X'$ flat over $S'$, $X = S \times_{S'} X'$, and
$Y = S \times_{S'} Y'$. Then
\begin{enumerate}
\item $f$ is flat if and only if $f'$ is flat,
\label{item-flat}
\item $f$ is an isomorphism if and only if $f'$ is an isomorphism,
\label{item-isomorphism}
\item $f$ is an open immersion if and only if $f'$ is an open immersion,
\label{item-open-immersion}
\item $f$ is quasi-compact if and only if $f'$ is quasi-compact,
\label{item-quasi-compact}
\item $f$ is universally closed if and only if $f'$ is universally closed,
\label{item-universally-closed}
\item $f$ is (quasi-)separated if and only if $f'$ is (quasi-)separated,
\label{item-separated}
\item $f$ is a monomorphism if and only if $f'$ is a monomorphism,
\label{item-monomorphism}
\item $f$ is surjective if and only if $f'$ is surjective,
\label{item-surjective}
\item $f$ is universally injective if and only if $f'$ is universally injective,
\label{item-universally-injective}
\item $f$ is affine if and only if $f'$ is affine,
\label{item-affine}
\item
\label{item-finite-type}
$f$ is locally of finite type if and only if $f'$ is locally of finite type,
\item $f$ is locally quasi-finite if and only if $f'$ is locally quasi-finite,
\label{item-quasi-finite}
\item
\label{item-finite-presentation}
$f$ is locally of finite presentation if and only if $f'$ is locally of
finite presentation,
\item
\label{item-relative-dimension-d}
$f$ is locally of finite type of relative dimension $d$ if and only if
$f'$ is locally of finite type of relative dimension $d$,
\item $f$ is universally open if and only if $f'$ is universally open,
\label{item-universally-open}
\item $f$ is syntomic if and only if $f'$ is syntomic,
\label{item-syntomic}
\item $f$ is smooth if and only if $f'$ is smooth,
\label{item-smooth}
\item $f$ is unramified if and only if $f'$ is unramified,
\label{item-unramified}
\item $f$ is \'etale if and only if $f'$ is \'etale,
\label{item-etale}
\item $f$ is proper if and only if $f'$ is proper,
\label{item-proper}
\item $f$ is integral if and only if $f'$ is integral,
\label{item-integral}
\item $f$ is finite if and only if $f'$ is finite,
\label{item-finite}
\item
\label{item-finite-locally-free}
$f$ is finite locally free (of rank $d$) if and only if $f'$
is finite locally free (of rank $d$), and
\item add more here.
\end{enumerate}
\end{lemma}

\begin{proof}
The assumptions on $X$ and $Y$ mean that $f$ is the base change of
$f'$ by $X \to X'$.
The properties $\mathcal{P}$ listed in (1) -- (23) above are all stable
under base change, hence if $f'$ has property $\mathcal{P}$, then so
does $f$. See
Schemes, Lemmas \ref{schemes-lemma-base-change-immersion},
\ref{schemes-lemma-quasi-compact-preserved-base-change},
\ref{schemes-lemma-separated-permanence}, and
\ref{schemes-lemma-base-change-monomorphism}
and
Morphisms, Lemmas
\ref{morphisms-lemma-base-change-surjective},
\ref{morphisms-lemma-base-change-universally-injective},
\ref{morphisms-lemma-base-change-affine},
\ref{morphisms-lemma-base-change-finite-type},
\ref{morphisms-lemma-base-change-quasi-finite},
\ref{morphisms-lemma-base-change-finite-presentation},
\ref{morphisms-lemma-base-change-relative-dimension-d},
\ref{morphisms-lemma-base-change-syntomic},
\ref{morphisms-lemma-base-change-smooth},
\ref{morphisms-lemma-base-change-unramified},
\ref{morphisms-lemma-base-change-etale},
\ref{morphisms-lemma-base-change-proper},
\ref{morphisms-lemma-base-change-finite}, and
\ref{morphisms-lemma-base-change-finite-locally-free}.

\medskip\noindent
The interesting direction in each case is therefore to assume
that $f$ has the property and deduce that $f'$ has it too.
By induction on the order of the thickening we may
assume that $S \subset S'$ is a first order thickening, see
discussion immediately following
Definition \ref{definition-thickening}.
We make a couple of general remarks which we will use without further
mention in the arguments below.
(I) Let $W' \subset S'$ be an affine open and let $U' \subset X'$
and $V' \subset Y'$ be affine opens lying over $W'$ with $f'(U') \subset V'$.
Let $W' = \Spec(R')$ and denote $I \subset R'$ be the ideal
defining the closed subscheme $W' \cap S$. Say $U' = \Spec(B')$
and $V' = \Spec(A')$. Then we get a commutative diagram
$$
\xymatrix{
0 \ar[r] &
IB' \ar[r] &
B' \ar[r] &
B \ar[r] & 0 \\
0 \ar[r] &
IA' \ar[r] \ar[u] &
A' \ar[r] \ar[u] &
A \ar[r] \ar[u] & 0
}
$$
with exact rows. Moreover $IB' \cong I \otimes_R B$, see proof of
Lemma \ref{lemma-deform}.
(II) The morphisms $X \to X'$ and $Y \to Y'$ are universal homeomorphisms.
Hence the topology of the maps $f$ and $f'$ (after any base change)
is identical. (III) If $f$ is flat, then $f'$ is flat and
$Y' \to S'$ is flat at every point in the image of $f'$, see
Lemma \ref{lemma-flatness-morphism-thickenings}.

\medskip\noindent
Ad (\ref{item-flat}). This is general remark (III).

\medskip\noindent
Ad (\ref{item-isomorphism}). Assume $f$ is an isomorphism.
By (III) we see that $Y' \to S'$ is flat. Choose an
affine open $V' \subset Y'$ and set $U' = (f')^{-1}(V')$. Then
$V = Y \cap V'$ is affine which implies that
$V \cong f^{-1}(V) = U = Y \times_{Y'} U'$ is affine. By
Lemma \ref{lemma-thickening-affine-scheme}
we see that $U'$ is affine. Thus we have a diagram as in the
general remark (I) and moreover $IA \cong I \otimes_R A$ because
$R' \to A'$ is flat. Then
$IB' \cong I \otimes_R B \cong I \otimes_R A \cong IA'$
and $A \cong B$. By the exactness of the rows
in the diagram above we see that $A' \cong B'$,
i.e., $U' \cong V'$. Thus $f'$ is an isomorphism.

\medskip\noindent
Ad (\ref{item-open-immersion}). Assume $f$ is an open immersion.
Then $f$ is an isomorphism of $X$ with an open subscheme $V \subset Y$.
Let $V' \subset Y'$ be the open subscheme whose underlying topological
space is $V$. Then $f'$ is a map from $X'$ to $V'$ which is an isomorphism by
(\ref{item-isomorphism}). Hence $f'$ is an open immersion.

\medskip\noindent
Ad (\ref{item-quasi-compact}). Immediate from remark (II). See also
Lemma \ref{lemma-thicken-property-morphisms} for a more general statement.

\medskip\noindent
Ad (\ref{item-universally-closed}). Immediate from remark (II). See also
Lemma \ref{lemma-thicken-property-morphisms} for a more general statement.

\medskip\noindent
Ad (\ref{item-separated}). Note that
$X \times_Y X = Y \times_{Y'} (X' \times_{Y'} X')$ so that
$X' \times_{Y'} X'$ is a thickening of $X \times_Y X$.
Hence the topology of the maps $\Delta_{X/Y}$ and $\Delta_{X'/Y'}$
matches and we win. See also
Lemma \ref{lemma-thicken-property-morphisms} for a more general statement.

\medskip\noindent
Ad (\ref{item-monomorphism}). Assume $f$ is a monomorphism.
Consider the diagonal morphism $\Delta_{X'/Y'} : X' \to X' \times_{Y'} X'$.
The base change of $\Delta_{X'/Y'}$ by $S \to S'$ is $\Delta_{X/Y}$
which is an isomorphism by assumption. By (\ref{item-isomorphism})
we conclude that $\Delta_{X'/Y'}$ is an isomorphism.

\medskip\noindent
Ad (\ref{item-surjective}). This is clear. See also
Lemma \ref{lemma-thicken-property-morphisms} for a more general statement.

\medskip\noindent
Ad (\ref{item-universally-injective}). Immediate from remark (II). See also
Lemma \ref{lemma-thicken-property-morphisms} for a more general statement.

\medskip\noindent
Ad (\ref{item-affine}). Assume $f$ is affine.  Choose an
affine open $V' \subset Y'$ and set $U' = (f')^{-1}(V')$.
Then $V = Y \cap V'$ is affine which implies that
$U = Y \times_{Y'} U'$ is affine. By
Lemma \ref{lemma-thickening-affine-scheme}
we see that $U'$ is affine. Hence $f'$ is affine. See also
Lemma \ref{lemma-thicken-property-morphisms} for a more general statement.

\medskip\noindent
Ad (\ref{item-finite-type}). Via remark (I) comes down to proving $A' \to B'$
is of finite type if $A \to B$ is of finite type. Suppose that
$x_1, \ldots, x_n \in B'$ are elements whose images in $B$ generate $B$
as an $A$-algebra. Then $A'[x_1, \ldots, x_n] \to B$ is surjective
as both $A'[x_1, \ldots, x_n] \to B$ is surjective and
$I \otimes_R A[x_1, \ldots, x_n] \to I \otimes_R B$ is surjective. See also
Lemma \ref{lemma-thicken-property-morphisms-cartesian}
for a more general statement.

\medskip\noindent
Ad (\ref{item-quasi-finite}). Follows from (\ref{item-finite-type}) and that
quasi-finiteness of a morphism of finite type can be checked on fibres, see
Morphisms, Lemma \ref{morphisms-lemma-quasi-finite-at-point-characterize}.
See also Lemma \ref{lemma-thicken-property-morphisms-cartesian}
for a more general statement.

\medskip\noindent
Ad (\ref{item-finite-presentation}). Via remark (I) comes down to proving
$A' \to B'$ is of finite presentation if $A \to B$ is of finite presentation.
We may assume that $B' = A'[x_1, \ldots, x_n]/K'$ for some ideal $K'$ by
(\ref{item-finite-type}). We get a short exact sequence
$$
0 \to K' \to A'[x_1, \ldots, x_n] \to B' \to 0
$$
As $B'$ is flat over $R'$ we see that $K' \otimes_{R'} R$ is the kernel of
the surjection $A[x_1, \ldots, x_n] \to B$. By assumption on $A \to B$ there
exist finitely many $f'_1, \ldots, f'_m \in K'$ whose images in
$A[x_1, \ldots, x_n]$ generate this kernel. Since $I$ is nilpotent we see
that $f'_1, \ldots, f'_m$ generate $K'$ by Nakayama's lemma, see
Algebra, Lemma \ref{algebra-lemma-NAK}.

\medskip\noindent
Ad (\ref{item-relative-dimension-d}). Follows from (\ref{item-finite-type})
and general remark (II). See also
Lemma \ref{lemma-thicken-property-morphisms-cartesian}
for a more general statement.

\medskip\noindent
Ad (\ref{item-universally-open}). Immediate from general remark (II). See also
Lemma \ref{lemma-thicken-property-morphisms} for a more general statement.

\medskip\noindent
Ad (\ref{item-syntomic}). Assume $f$ is syntomic. By
(\ref{item-finite-presentation}) $f'$ is locally of finite presentation,
by general remark (III) $f'$ is flat and the fibres of $f'$ are the fibres
of $f$. Hence $f'$ is syntomic by
Morphisms, Lemma \ref{morphisms-lemma-syntomic-flat-fibres}.

\medskip\noindent
Ad (\ref{item-smooth}). Assume $f$ is smooth. By
(\ref{item-finite-presentation}) $f'$ is locally of finite presentation,
by general remark (III) $f'$ is flat, and the fibres of $f'$ are the
fibres of $f$. Hence $f'$ is smooth by
Morphisms, Lemma \ref{morphisms-lemma-smooth-flat-smooth-fibres}.

\medskip\noindent
Ad (\ref{item-unramified}). Assume $f$ unramified. By
(\ref{item-finite-type}) $f'$ is locally of finite type
and the fibres of $f'$ are the fibres of $f$.
Hence $f'$ is unramified by
Morphisms, Lemma \ref{morphisms-lemma-unramified-etale-fibres}. See also
Lemma \ref{lemma-thicken-property-morphisms-cartesian}
for a more general statement.

\medskip\noindent
Ad (\ref{item-etale}). Assume $f$ \'etale. By
(\ref{item-finite-presentation}) $f'$ is locally of finite presentation,
by general remark (III) $f'$ is flat, and the fibres of $f'$ are the fibres
of $f$. Hence $f'$ is \'etale by
Morphisms, Lemma \ref{morphisms-lemma-etale-flat-etale-fibres}.

\medskip\noindent
Ad (\ref{item-proper}). This follows from a combination of
(\ref{item-separated}), (\ref{item-finite-type}), (\ref{item-quasi-compact}),
and (\ref{item-universally-closed}).  See also
Lemma \ref{lemma-thicken-property-morphisms-cartesian}
for a more general statement.

\medskip\noindent
Ad (\ref{item-integral}). Combine (\ref{item-universally-closed}) and
(\ref{item-affine}) with
Morphisms, Lemma \ref{morphisms-lemma-integral-universally-closed}. See also
Lemma \ref{lemma-thicken-property-morphisms} for a more general statement.

\medskip\noindent
Ad (\ref{item-finite}). Combine (\ref{item-integral}),
and (\ref{item-finite-type}) with
Morphisms, Lemma \ref{morphisms-lemma-finite-integral}. See also
Lemma \ref{lemma-thicken-property-morphisms-cartesian}
for a more general statement.

\medskip\noindent
Ad (\ref{item-finite-locally-free}). Assume $f$ finite locally free. By
(\ref{item-finite}) we see that $f'$ is finite, by general remark (III)
$f'$ is flat, and by (\ref{item-finite-presentation}) $f'$ is locally of finite
presentation. Hence $f'$ is finite locally free by
Morphisms, Lemma \ref{morphisms-lemma-finite-flat}.
\end{proof}

\end{document}
